\section{Numerical Prediction}

\subsection{From fixed point to polynomial}

The previous section derived the second-order compact-fiber bridge
equation
\[
x
=
137+\frac{74}{15x}
-
\frac{296}{105x^3}.
\]
Multiplying by \(x^3\) converts this fixed-point equation into a quartic
polynomial whose positive root defines the compact-fiber inverse bridge
coupling:
\[
\boxed{
x^4
-
137x^3
-
\frac{74}{15}x^2
+
\frac{296}{105}
=
0.
}
\]
The measured value of \(\alpha\) is not used to obtain the equation or
its root.

\subsection{Positive root}

Solving
\[
x^4
-
137x^3
-
\frac{74}{15}x^2
+
\frac{296}{105}
=
0
\]
gives the physically relevant positive root
\[
\boxed{
x_{\rm fiber}=137.0359991771859\ldots.
}
\]
Therefore the compact-fiber electromagnetic bridge coupling is
\[
\boxed{
\alpha_F
=
\frac1{x_{\rm fiber}}.
}
\]
This is the central numerical output of the compact-fiber bridge
equation.

\subsection{First-order and second-order values}

The first-order equation
\[
x=B+\frac{K}{x},
\qquad
B=137,
\qquad
K=\frac{74}{15},
\]
gives
\[
\boxed{
x_{\rm quad}=137.03600027236\ldots.
}
\]
The second-order closed-return equation replaces \(x_{\rm quad}\) with
the positive quartic root \(x_{\rm fiber}\). The shift is produced by
the admitted scalar two-node return term
\[
-\frac{296}{105x^3}.
\]
The first-order result fixes the inverse bridge value close to the
empirical low-energy inverse fine-structure constant scale but remains
high by approximately
\[
8.0\ \text{ppb}.
\]
The second-order closed return supplies the refinement admitted by the
framework rules stated above.

\subsection{Comparison with the empirical low-energy value}

Using the CODATA 2022 recommended value
\cite{NIST_CODATA_2022_alpha_inv,NIST_CODATA_2022_wallet}
\[
\alpha^{-1}_{\rm CODATA}=137.035999177(21),
\]
the compact-fiber value
\[
x_{\rm fiber}=137.0359991771859\ldots
\]
lies within the experimental uncertainty band of the quoted value. The
difference from the CODATA 2022 central value is approximately
\[
1.4\times10^{-3}\ \text{ppb},
\]
while the quoted experimental uncertainty is approximately
\[
0.15\ \text{ppb}.
\]

The numerical comparison is not used to construct the bridge equation.
It is also not a statistical validation of the framework, since the
present paper does not assign a theoretical uncertainty to the imported
bridge-readout rules, conditional framework lemmas, or possible omitted
same-order mechanisms. The comparison is therefore reported as a
consistency check between the compact-fiber point value and the empirical
low-energy coupling, not as an error-budget closure.

\subsection{Status of the prediction}

No measured value of \(\alpha\) appears in the compact-fiber bridge
equation, and no continuously adjusted coefficient appears in the
polynomial. The coefficients are fixed structurally before comparison
with CODATA:
\[
137=128+9,
\qquad
\frac{74}{15}
=
5-\frac1{15},
\qquad
\frac{296}{105}
=
\left(\frac{74}{15}\right)\left(\frac47\right).
\]
Their structural origins are
\[
128=|Z_4\times Z_4\times Z_8|,
\qquad
9=|V_{\rm rad}^{\rm bridge}|=3^2,
\]
\[
5=N_m+1,
\qquad
\frac1{15}=\frac1{N_m^2-1},
\qquad
\frac47=\frac{N_m}{N_e-1}.
\]

Given the stated compact-fiber bridge lemmas, the quartic equation and
its positive root follow algebraically. The result remains conditional
on the bridge-readout, closure-budget, and closed-return rules stated
above; the paper does not claim that those rules are derived here from
bare postulates alone. The role identification
\[
\alpha_F=\alpha_{\rm low-energy}
\]
is supported by the numerical agreement, but is not based on numerical
agreement alone. Its structural basis is the material/radiative bridge
role of \(\alpha_F\) and the two-vertex \(\alpha_F^2\) dependence of
leading atomic binding.